\documentclass{article}
\usepackage{graphicx}
\usepackage[margin=1.5cm]{geometry}
\usepackage{amsmath}

\begin{document}
\twocolumn

\title{Wednesday Warm Up: Unit 7: Forces V and Uniform Circular Motion, II}
\author{Prof. Jordan C. Hanson}

\maketitle

\section{Memory Bank}

\begin{itemize}
\item $\tau = \vec{r} \times \vec{F}$ ... Definition of torque.
\item $W = \vec{F} \cdot \Delta \vec{x}$ ... Definition of work (linear).
\item $W = \vec{\tau} \cdot \Delta \vec{\theta}$ ... Definition of work (angular).
\item $\tau = I\alpha$ ... Magnitude of work, angular acceleration, and moment of inertia, $I$.
\item $I = mr^2$ ... Moment of inertia of a point mass $m$ rotating with radius $r$.
\item \textbf{For sytems in static equilibrium, (1) the net force is zero, and (2) the net torque is zero.}
\item $\vec{L} = \vec{r} \times \vec{p}$ ... Definition of \textit{angular momentum}.
\item $L = I\omega$, and $\vec{L} = I\vec{\omega}$, and $\Delta \vec{L} / \Delta t = I\vec{\alpha}$
\end{itemize}

\section{Dot Product, and Cross Product}

\begin{enumerate}
\item Find the \textit{angular momentum}:
\begin{itemize}
\item $\vec{r} = -2\hat{i} + \hat{j}$ m, $\vec{p} = 4\hat{i} + \hat{j}$ kg m s$^{-1}$.
\item $\vec{r} = 2\hat{i} + 1\hat{j}$ m, $\vec{p} = 1\hat{i} - 2\hat{j}$ kg m s$^{-1}$.
\end{itemize}
\vspace{1.5cm}
\end{enumerate}

\section{Torque and Angular Momentum}

\begin{enumerate}
\item A motorcycle wheel has a mass of 12.0 kg. Assume it is a disk of radius 0.330 m. ($I = \frac{1}{2} M R^2$).  The motorcycle is on its center stand, so that the wheel can spin freely. (a) If the drive chain exerts a force of 3000 N at a radius of 6.00 cm, what is the angular acceleration of the wheel? (b) What is the tangential acceleration of a point on the outer edge of the tire? (c) How long does it take to reach 80.0 rad/s, starting from rest? What is this in rpm? \\ \vspace{3cm}
\item Once the wheel has reached 80 rad/s, what is the angular momentum? \\ \vspace{2cm}
\item Suppose a loose part of the wheel fell off while the system is spinning at 80 rad/s.  If the mass drops by 20\%, what is the new angular velocity? \\ \vspace{2cm}
\item Remember the example of \textit{las bolas}, the two balls spinning around a common radius, connected by a chord.  (a) Show that the angular momentum is $L = 2 mr^2 \omega$, where $m$ is the mass of one \textit{bola}, $r$ is the radius, and $\omega$ is the angular velocity.
\end{enumerate}

\end{document}
